\chapter{总体描述}

本章从宏观角度描述电力数据分析系统的产品定位、功能概述、目标用户、约束条件和外部依赖，为后续详细需求描述提供整体框架。

\section{产品前景}

\subsection{项目背景}

随着可再生能源的普及和智能电网技术的发展，家庭和商业用户越来越需要智能化的能源管理解决方案。分时电价（Time-of-Use, TOU）政策的推广使得用户可以通过优化用电时段来降低电费支出。储能电池技术的成熟为实现"低价充电、高价放电"的能源套利策略提供了硬件基础。

本系统正是在这一背景下设计开发的，旨在利用机器学习技术预测能源负载，结合数学优化算法生成最优电池调度策略，帮助用户实现能源成本最小化。

\subsection{系统定位}

电力数据分析系统是一个独立运行的云端应用，但需要与多个外部系统进行数据交互，如图~\ref{fig:system-context} 所示。

\begin{figure}[htbp]
    \centering
    \begin{tikzpicture}[
            node distance=2.5cm,
            system/.style={rectangle, draw=blue!70, fill=blue!15,
                    text width=4cm, minimum height=2cm,
                    align=center, rounded corners=5pt,
                    font=\small\bfseries, line width=1.5pt},
            external/.style={rectangle, draw=gray!60, fill=gray!10,
                    text width=3cm, minimum height=1.5cm,
                    align=center, rounded corners=3pt,
                    font=\footnotesize},
            user/.style={ellipse, draw=green!60, fill=green!10,
                    minimum width=2.5cm, minimum height=1.5cm,
                    align=center, font=\footnotesize},
            arrow/.style={->, >=stealth, thick},
            darrow/.style={<->, >=stealth, thick}
        ]
        % 核心系统
        \node[system] (core) {电力数据分析系统\\{\footnotesize (Power Data Analysis System)}};

        % 外部系统
        \node[external, above left=2cm and 3cm of core] (caiso) {CAISO\\电网负载数据};
        \node[external, above=2.5cm of core] (weather) {OpenWeatherMap\\天气数据API};
        \node[external, above right=2cm and 3cm of core] (firebase) {Firebase\\认证服务};

        \node[external, below left=2cm and 3cm of core] (gcs) {Google Cloud\\Storage};
        \node[external, below=2.5cm of core] (firestore) {Cloud\\Firestore};
        \node[external, below right=2cm and 3cm of core] (gae) {Google App\\Engine};

        % 用户
        \node[user, left=4cm of core] (user) {终端用户\\(Web浏览器)};

        % 连接线
        \draw[arrow] (caiso) -- node[above, sloped, font=\tiny] {负载数据} (core);
        \draw[arrow] (weather) -- node[right, font=\tiny] {温度预报} (core);
        \draw[darrow] (firebase) -- node[above, sloped, font=\tiny] {认证令牌} (core);

        \draw[darrow] (core) -- node[above, sloped, font=\tiny] {文件存储} (gcs);
        \draw[darrow] (core) -- node[left, font=\tiny] {历史记录} (firestore);
        \draw[arrow] (gae) -- node[above, sloped, font=\tiny] {托管运行} (core);

        \draw[darrow] (user) -- node[above, font=\tiny] {HTTP/HTTPS} (core);
    \end{tikzpicture}
    \caption{系统上下文图}
    \label{fig:system-context}
\end{figure}

\subsection{与其他系统的关系}

表~\ref{tab:external-systems} 描述了本系统与外部系统的交互关系。

\begin{table}[htbp]
    \centering
    \caption{外部系统交互关系}
    \label{tab:external-systems}
    \begin{tabularx}{\textwidth}{@{}l l X@{}}
        \toprule
        \textbf{外部系统}           & \textbf{交互方式} & \textbf{交互内容}         \\
        \midrule
        CAISO API               & HTTP GET      & 获取加州电网实时负载数据，用于预测模型训练 \\
        OpenWeatherMap API      & HTTP GET      & 获取城市温度预报，作为预测特征       \\
        Firebase Authentication & OAuth 2.0     & 用户身份认证和会话管理           \\
        Google Cloud Storage    & REST API      & 存储用户数据文件和机器学习模型       \\
        Cloud Firestore         & SDK/REST      & 存储用户历史记录和系统配置         \\
        Google App Engine       & PaaS 托管       & 运行后端 Flask 应用和定时任务    \\
        \bottomrule
    \end{tabularx}
\end{table}

\section{产品功能概述}

系统的主要功能模块如图~\ref{fig:function-overview} 所示，可分为四大类：

\begin{figure}[H]
    \centering
    \begin{tikzpicture}[
            node distance=0.5cm,
            category/.style={rectangle, draw=blue!60, fill=blue!20,
                    text width=4cm, minimum height=1cm,
                    align=center, rounded corners=3pt,
                    font=\small\bfseries},
            feature/.style={rectangle, draw=gray!50, fill=white,
                    text width=3.8cm, minimum height=0.7cm,
                    align=left, font=\footnotesize,
                    left color=white, right color=gray!10}
        ]
        % 用户认证模块
        \node[category] (auth) {用户认证模块};
        \node[feature, below=0.3cm of auth] (auth1) {• Email/密码登录};
        \node[feature, below=0.1cm of auth1] (auth2) {• Google OAuth登录};
        \node[feature, below=0.1cm of auth2] (auth3) {• 会话管理与登出};

        % 核心业务模块
        \node[category, right=2cm of auth] (core) {核心业务模块};
        \node[feature, below=0.3cm of core] (core1) {• 能源负载预测 (ML)};
        \node[feature, below=0.1cm of core1] (core2) {• 电池调度优化 (Gurobi)};
        \node[feature, below=0.1cm of core2] (core3) {• 实时数据采集};
        \node[feature, below=0.1cm of core3] (core4) {• What-If 场景分析};

        % 数据管理模块
        \node[category, below=3cm of auth] (data) {数据管理模块};
        \node[feature, below=0.3cm of data] (data1) {• 文件上传/下载};
        \node[feature, below=0.1cm of data1] (data2) {• 历史记录查询};
        \node[feature, below=0.1cm of data2] (data3) {• 数据质量分析};

        % 可视化模块
        \node[category, right=2cm of data] (viz) {数据可视化模块};
        \node[feature, below=0.3cm of viz] (viz1) {• 负载/价格曲线图};
        \node[feature, below=0.1cm of viz1] (viz2) {• 电池SOC变化图};
        \node[feature, below=0.1cm of viz2] (viz3) {• 特征重要性图表};
        \node[feature, below=0.1cm of viz3] (viz4) {• 成本节省统计};
    \end{tikzpicture}
    \caption{产品功能概述图}
    \label{fig:function-overview}
\end{figure}

\subsection{功能列表}

\begin{enumerate}
    \item \textbf{用户认证与权限管理}
          \begin{itemize}
              \item 支持 Email/密码注册和登录
              \item 支持 Google 账户第三方登录
              \item 自动会话保持和安全登出
          \end{itemize}

    \item \textbf{能源负载预测}
          \begin{itemize}
              \item 基于随机森林/集成学习的 24 小时负载预测
              \item 支持温度调整的 What-If 分析
              \item 模型可解释性展示（SHAP 特征贡献）
          \end{itemize}

    \item \textbf{电池调度优化}
          \begin{itemize}
              \item 基于 Gurobi 求解器的线性规划优化
              \item 生成最优充放电时刻表
              \item 计算成本节省金额和百分比
          \end{itemize}

    \item \textbf{数据管理}
          \begin{itemize}
              \item CSV 文件上传和分析
              \item 历史分析记录查询
              \item 实时数据自动采集和存储
          \end{itemize}

    \item \textbf{数据可视化}
          \begin{itemize}
              \item 交互式图表展示（基于 fl\_chart）
              \item 响应式仪表盘布局
              \item 关键指标卡片展示
          \end{itemize}
\end{enumerate}

\section{用户特征}

\subsection{目标用户群体}

本系统的目标用户主要包括以下群体：

\begin{table}[htbp]
    \centering
    \caption{用户特征描述}
    \label{tab:user-characteristics}
    \begin{tabularx}{\textwidth}{@{}l l X@{}}
        \toprule
        \textbf{用户类型} & \textbf{技术背景} & \textbf{使用场景}      \\
        \midrule
        家庭用户          & 初级            & 查看能源优化建议，了解潜在电费节省  \\
        能源管理专员        & 中级            & 配置电池参数，分析多种优化场景    \\
        数据科学学生        & 中高级           & 学习 ML 预测和优化算法的实际应用 \\
        系统管理员         & 高级            & 监控系统运行状态，管理数据和模型   \\
        \bottomrule
    \end{tabularx}
\end{table}

\subsection{用户能力假设}

\begin{itemize}
    \item 用户具备基本的计算机操作能力，能够使用 Web 浏览器
    \item 用户了解基本的能源概念（如电价、用电量、储能电池）
    \item 高级用户可能需要理解机器学习和优化算法的基本原理
\end{itemize}

\section{一般约束}

\subsection{技术约束}

\begin{enumerate}
    \item \textbf{云平台依赖}：系统部署于 Google Cloud Platform，受限于 GCP 服务的可用性和地理覆盖。

    \item \textbf{免费层额度限制}：
          \begin{itemize}
              \item Google App Engine：每日 28 实例小时（F1 实例）
              \item Cloud Storage：5 GB 存储，1 GB 出站流量/月
              \item Firestore：1 GiB 存储，50,000 次读取/日
          \end{itemize}

    \item \textbf{Gurobi 许可证}：优化求解器需要有效的 Gurobi 许可证（学术版免费）。

    \item \textbf{第三方 API 限制}：OpenWeatherMap 免费版限制 60 次调用/分钟。
\end{enumerate}

\subsection{开发约束}

\begin{enumerate}
    \item \textbf{开发语言}：后端使用 Python $\geq$ 3.9，前端使用 Dart (Flutter Web)
    \item \textbf{开发周期}：项目需在规定的学期内完成开发和测试
    \item \textbf{团队规模}：开发团队由 3 名成员组成
\end{enumerate}

\subsection{运营约束}

\begin{enumerate}
    \item 系统运行需要稳定的互联网连接
    \item 用户浏览器需支持现代 Web 标准（Chrome、Firefox、Safari、Edge）
\end{enumerate}

\section{假设与依赖}

\subsection{假设条件}

\begin{enumerate}
    \item 假设用户拥有稳定的互联网连接
    \item 假设 Google Cloud Platform 服务保持正常可用
    \item 假设 CAISO 和 OpenWeatherMap API 持续提供数据服务
    \item 假设用户输入的电池参数在合理范围内
    \item 假设分时电价政策在预测周期内保持稳定
\end{enumerate}

\subsection{外部依赖}

表~\ref{tab:dependencies} 列出了系统的主要外部依赖项。

\begin{table}[htbp]
    \centering
    \caption{外部依赖项}
    \label{tab:dependencies}
    \begin{tabularx}{\textwidth}{@{}l l X@{}}
        \toprule
        \textbf{依赖项} & \textbf{版本/来源} & \textbf{用途} \\
        \midrule
        Python       & $\geq$ 3.9     & 后端运行时环境     \\
        Flask        & $\geq$ 2.0     & Web 框架      \\
        Pandas       & $\geq$ 1.5     & 数据处理        \\
        Scikit-learn & $\geq$ 1.0     & 机器学习模型      \\
        Gurobi       & $\geq$ 10.0    & 优化求解器       \\
        Flutter      & $\geq$ 3.0     & 前端框架        \\
        Firebase SDK & 最新稳定版          & 认证和托管服务     \\
        \bottomrule
    \end{tabularx}
\end{table}

\section{系统架构概览}

图~\ref{fig:system-architecture} 展示了电力数据分析系统的三层架构设计。

\begin{figure}[htbp]
    \centering
    \includegraphics[width=0.85\textwidth]{GAE.png}
    \caption{系统部署架构图（Google App Engine）}
    \label{fig:system-architecture}
\end{figure}

系统采用前后端分离的三层架构：

\begin{enumerate}
    \item \textbf{表示层（Presentation Layer）}：基于 Flutter Web 开发，部署于 Firebase Hosting，负责用户界面展示和交互。

    \item \textbf{业务逻辑层（Business Logic Layer）}：基于 Python Flask 开发，部署于 Google App Engine，负责 ML 预测、优化计算和 API 服务。

    \item \textbf{数据层（Data Layer）}：整合 Google Cloud Storage（文件存储）和 Cloud Firestore（结构化数据），负责数据持久化。
\end{enumerate}

图~\ref{fig:data-flow} 展示了系统的数据流转过程。

\begin{figure}[H]
    \centering
    \includegraphics[width=0.6\textwidth]{作业图.png}
    \caption{系统数据流程图}
    \label{fig:data-flow}
\end{figure}